\section{Results and Analysis}
\label{sec:results}
\subsection{Expected Value of Perfect Information}

The Expected Value of Perfect Information (EVPI) quantifies the expected gain from knowing
the exact scenario before making any decisions. It is estimated as follows.

First, the stochastic MILP is solved using $S$ training scenarios drawn from the
chosen sampling distribution, yielding a here-and-now policy
$(X^*, Y^*, A^*, D^*)$.
Second, $I$ independent evaluation scenarios are generated.
For each evaluation scenario $i$, two objective values are computed:
\begin{itemize}
    \item $z_{\text{policy}}(i)$: the cost using the policy trained on S evaluated on scenario $i$.
    \item $z_{\text{perf}}(i)$: the cost with perfect information, obtained by solving the MILP with scenario $i$ as the sole scenario, so all first-stage decisions can be optimised with full knowledge of the realised durations.
\end{itemize}
The value of perfect information for scenario $i$ is
\[
    \text{VPI}(i) = z_{\text{policy}}(i) - z_{\text{perf}}(i) \;\geq\; 0,
\]
and the EVPI is estimated as the sample mean over the $I$ evaluation scenarios:
\[
    \widehat{\text{EVPI}} = \frac{1}{I} \sum_{i=1}^{I} \text{VPI}(i).
\]

Table~\ref{tab:vpi} presents the per-scenario costs for 10 training scenarios
and 10 evaluation scenarios. In 8 out of 10 evaluation scenarios,
$z_{\text{perf}}(i) = 0$, meaning that a penalty-free schedule exists when
the surgery durations are known in advance. For these scenarios, the entire
cost incurred by the policy is attributable to uncertainty, and
$\text{VPI}(i) = z_{\text{policy}}(i)$. Only 2 scenarios (i = 1 and i = 4)
require penalties even under perfect information, suggesting that these
particular realisations of surgery durations are inherently infeasible to
schedule without cost regardless of the decision rule applied. The estimated
expected value of perfect information over the 10 evaluation scenarios is
$\widehat{\text{EVPI}} = 158.05$, indicating that on average a large share of
the policy cost could in principle be eliminated by resolving duration
uncertainty prior to scheduling.

\begin{table}[h]
    \centering
    \caption{Per-scenario costs and Value of Perfect Information (VPI).}
    \label{tab:vpi}
    \begin{tabular}{rrrr}
        \toprule
        $i$ & $z_{\text{policy}}(i)$ & $z_{\text{perf}}(i)$ & $\text{VPI}(i)$ \\
        \midrule
         1 & 225.146 & 16.048 & 209.098 \\
         2 & 119.461 &  0.000 & 119.461 \\
         3 & 119.849 &  0.000 & 119.849 \\
         4 & 177.055 &  2.742 & 174.313 \\
         5 & 153.461 &  0.000 & 153.461 \\
         6 & 128.878 &  0.000 & 128.878 \\
         7 & 140.100 &  0.000 & 140.100 \\
         8 & 186.231 &  0.000 & 186.231 \\
         9 & 174.386 &  0.000 & 174.386 \\
        10 & 174.679 &  0.000 & 174.679 \\
        \bottomrule
    \end{tabular}
\end{table}